\casdocAbsLabel{Objective:}
This public example demonstrates a complete postdoctoral research report template for the Chinese Academy of Sciences, using the image circulation and aesthetic reception of Nozomi Sasaki in East Asian popular culture as the sample topic.

\casdocAbsLabel{Methods:}
The demonstration adopts a five-chapter structure covering research background, material organization, circulation paths, reception mechanisms, and conclusions. At the engineering level, it reuses the stable \texttt{bensz-thesis} toolchain while registering an independent \texttt{thesis-cas-postdoc} template, profile, and style so the project remains self-contained when distributed on its own.

\casdocAbsLabel{Results:}
The project delivers a fully compilable report scaffold whose cover reflects the key fields required by the national postdoctoral report guideline, including classification number, UDC, security level, document number, completion period, and submission date. The title page further records the postdoctoral researcher, discipline station, specialty, work period, and affiliation.

\casdocAbsLabel{Conclusion:}
The template shows that the existing thesis infrastructure can support the postdoctoral report scenario with an independent template identity, while keeping the build workflow stable and without affecting other thesis templates in the repository.

\casdocAbsLabel{Keywords:}
postdoctoral research report; Chinese Academy of Sciences; Nozomi Sasaki; East Asian popular culture; LaTeX template
